\section{Architecture}
\label{sec:arch}

The two agents share one BDI pipeline; only the top of the stack differs
(Agent~B adds the LLM and team layers). Each server event triggers the loop:

\begin{verbatim}
  socket events (map, config, you, sensing, msg)
            |
        BeliefBase        map graph, me, parcels, agents, mission, claims
            |
   OptionGenerator        what is possible (pickup / deliver / explore / wait)
            |
       Strategy           what is preferable  (swappable; LLM adapts it via
            |              mission state for Agent B)
   IntentionRevision      commit / revise (replace policy + hysteresis)
            |
     PlanLibrary          how: GoPickUp, DeliverCarried, Explore, Wait,
            |              mission + handover plans (Agent B),
            |              PddlGoTo -> FollowPathGoTo (BFS fallback)
     ActionExecutor       serialized move/pickup/putdown, red-light gate
\end{verbatim}

\paragraph{Control loop.} On each \texttt{sensing}/\texttt{you} event (and a
periodic timer, since sensing is change-driven) the agent revises beliefs,
regenerates options, and lets intention revision decide whether to keep the
current intention or switch. A new option must beat the current one by a
hysteresis margin before the agent commits, which prevents target thrashing.
The chosen intention is achieved by the first applicable plan in the plan
library; failures fall through to the next plan.

\paragraph{Design principles.} (i)~\emph{Infrastructure vs.\ strategy}:
strategies only score options; they never move the agent, so a new strategy is
one file (Section~\ref{sec:bdi}). (ii)~\emph{Deterministic low-level
movement}: pathfinding is BFS on a directed graph; the LLM stays high-level and
never drives motion. (iii)~\emph{Serialized actions}: all actuator calls are
chained into a promise queue to respect the server's action mutex.

\paragraph{Repository layout.}
\begin{center}\small
\begin{tabular}{@{}ll@{}}
\toprule
Path & Role \\
\midrule
\texttt{src/core/} & BeliefBase, OptionGenerator, IntentionRevision, Intention, PlanLibrary, ActionExecutor \\
\texttt{src/planning/} & GridGraph (directed, arrow tiles), PathPlanner (BFS), PddlPlanner \\
\texttt{src/strategies/} & four baseline strategies $+$ a registry (Section~\ref{sec:bdi}) \\
\texttt{src/llm/} & mission interpreter + tools (Agent~B, Section~\ref{sec:llm}) \\
\texttt{src/communication/} & team protocol (Section~\ref{sec:coord}) \\
\texttt{src/metrics/} & run logging and result aggregation \\
\texttt{scripts/} & run / experiment / baseline / campaign / smoke-test \\
\texttt{experiments/} & results and the committed baseline write-up \\
\bottomrule
\end{tabular}
\end{center}

\paragraph{Verification tooling.} An offline smoke test (\texttt{npm test}, no
server or network) pins the load-bearing semantics: directed-graph
construction with one-way arrow tiles, BFS and reversed multi-source delivery
distances, belief revision (whole-tick decay projection, negative evidence),
strategy selection, mission parsing, PDDL problem generation, and ack
normalization. It runs in milliseconds and gates every change.

\paragraph{Robustness note.} We never trust a single server ack shape: the
running course servers differ from the SDK type declarations. Pickup/putdown
acks may carry \texttt{\{id\}} objects, plain strings, or no usable id at all
(observed both locally and on the course server). All ack ids pass through a
normalization layer, with belief-reconciliation fallbacks when the server
contradicts the carry belief --- a fix prompted by a live ``phantom-carry''
loop (Section~\ref{sec:bdi}).
